\section{Introduction}
\label{sec:intro}
Counterfeit detection is a natural use for a cheap classifier. A scanner photographs a note, a few image statistics are computed, and a model flags it as genuine or forged. The UCI Banknote Authentication data \cite{lohweg2012uci} is a standard small benchmark for this task, with four wavelet features per image. On it, standard methods are almost perfect: a tuned logistic regression misclassifies \LRcv{} of notes in 5-fold cross-validation (CV). Ranking methods by their final accuracy here would mostly compare noise, because with a few hundred held-out notes the difference between 0.3\% and 0.6\% error is one or two notes.

We therefore ask a question that still separates methods when accuracy saturates: \emph{how many labelled training notes does each method need to reach a given error rate?} In practice, labelled forgeries are the expensive part of building such a system, because each one must be produced or seized and then verified. A method that reaches 1\% error from 100 notes is more useful than one that needs 800, even if both are perfect at 1,000.

We compare two methods taught early in the course, logistic regression (the baseline) and $k$-nearest neighbours ($k$NN), and three later ones as an extension: random forests, gradient boosting and a one-hidden-layer neural network. We draw learning curves from 10 to 800 training notes with 50 random draws per size, tune every method by CV inside each draw, and compare methods on the same draws. Our contributions are:
\begin{itemize}
  \item a leakage-aware protocol for this dataset: duplicate removal, scaling and tuning inside each training set, an untouched test set, and a check for several images of the same physical specimen;
  \item paired learning curves with standard errors, and the number of notes each method needs to reach 2\% and 1\% error, with bootstrap intervals;
  \item an honest reading of where the differences are real and where they are ties, including the one-time test-set result.
\end{itemize}
Among the two core methods we find a crossover. Logistic regression learns fastest at first, reaching 2\% error from about \reachLRTwo{} notes, but levels off near 0.6\% because no linear boundary separates the classes. $k$NN starts worse, catches up at about 100 notes and reaches 1\% from about \reachKNNOne{} notes, against \reachLROne{} for logistic regression. Among the extensions, the tree ensembles are the least sample-efficient methods we tried, while a small neural network with 8 hidden units has the lowest mean error at every size from 10 notes up and reaches 1\% from about \reachNNOne{} notes. On the 270 test notes, all five methods make between 0 and 3 errors, so the test set cannot separate them; the learning curves can.

\section{Related work}
\label{sec:related}
\paragraph{Dataset.} The data come from the authentication work of Gillich and Lohweg \cite{gillich2010banknote}: images of genuine and forged banknote-like specimens were taken with an industrial camera and summarised by statistics of a wavelet transform. Lohweg released the four features on the UCI repository \cite{lohweg2012uci}. Because most standard methods score close to 100\% on it, a single accuracy figure says little about which method to prefer.

\paragraph{Learning curves for comparing classifiers.} Perlich, Provost and Simonoff \cite{perlich2003tree} compared logistic regression and tree induction across training sizes in a large-scale study of many datasets. They found that logistic regression tends to win on small samples and trees on large ones, so that the curves can cross and a comparison at one size can mislead. Ng and Jordan \cite{ng2002discriminative} showed a related effect between generative and discriminative models: the method with the lower asymptotic error is not always the faster learner. Viering and Loog \cite{viering2023shape} review the shapes learning curves take, including non-monotone ones. Our study applies this learning-curve view to a dataset where final accuracy is saturated.

\paragraph{Fair evaluation.} Choosing hyperparameters and estimating error on the same folds biases the estimate downwards \cite{varma2006bias,cawley2010overfitting}, so we use nested CV for the baseline and tune inside every training draw. Leakage between training and evaluation data, for example near-copies of the same object, can make a model look better than it is \cite{kaufman2012leakage}. To compare two classifiers on one test set we use McNemar's test, as recommended by Dietterich \cite{dietterich1998approximate}.

\paragraph{Methods.} We use standard forms of logistic regression and CV \cite{james2021isl}, nearest-neighbour classification \cite{cover1967nearest}, random forests \cite{breiman2001random}, gradient boosting \cite{friedman2001greedy} and a small multilayer perceptron, all in scikit-learn \cite{pedregosa2011scikit}.
